\section{Derivation of Expected Precision Goal Stop Iteration}\label{app:expected_stop}

The precision-based stopping algorithms (PitG \citealt{kruschke2015doing} and DPitG) halt when the HDI width
(see Appendix~\ref{app:hdi})
falls below the precision goal $\omega_{\rm goal}$. Because the posterior HDI width
shrinks as $N^{-1/2}$, one can predict in advance the iteration at which this
threshold will be crossed. The General CLT Form below uses a normal approximation
and is intended as an analytical guide; the Binomial Case refines this using the
exact Beta posterior variance, and the Continuous Case provides two variants
(CLT and Student-$t$).
% Throughout, we assume unimodal, approximately symmetric posteriors, for which
% HDI width is a natural and interpretable precision measure.
% For multimodal posteriors, HDI width becomes ambiguous (it is unclear which mode
% to centre on), and posterior entropy may be a more appropriate precision criterion
% \citep{kruschke2015doing, chaloner1995, sebastiani2000}; such settings fall outside
% the scope of the present work.

\subsection*{General CLT Form}

Assume an estimator $\hat\theta_N$ with asymptotic distribution
\begin{equation}\label{eq:clt_general}
\hat\theta_N \approx \mathcal{N}\!\left(\theta,\,\frac{V(\theta)}{N}\right),
\end{equation}
where $V(\theta)$ is the per-observation variance contribution. For a symmetric
central interval with critical value $z_{*}$ (e.g., $z_{*}=1.96$ for 95\%), the
interval width is approximately
\begin{equation}\label{eq:hdi_width_approx}
w \approx 2 z_{*}\sqrt{\frac{V(\theta)}{N}}.
\end{equation}
Setting $w = \omega_{\rm goal}$ and solving for $N$ gives the general expected stop
iteration:
\begin{equation}\label{eq:expected_stop_iteration_clt}
N_{\rm goal} \approx \frac{4 z_{*}^{2}}{\omega_{\rm goal}^{2}}\,V(\theta).
\end{equation}

\subsection*{Binomial Case}

For Bernoulli outcomes the posterior under the Haldane prior $\mathrm{Beta}(0,0)$
is $\mathrm{Beta}(s,\,f)$, where $s$ is the number of successes and
$f = N - s$ the number of failures.\footnote{This appendix uses the
Haldane prior rather than the flat $\mathrm{Beta}(1,1)$ prior of
Appendix~\ref{app:hdi}: the Haldane convention yields the clean $-1$ correction
in the planning formula and is consistent with the planning function in the
accompanying code (\texttt{binomial\_rate\_ci\_width\_to\_sample\_size} in
\texttt{utils\_stats.py}).}
Its exact variance is
\[
\mathrm{Var}(\theta \mid s, f) = \frac{\hat\theta\,(1-\hat\theta)}{N+1},
\]
where $\hat\theta = s/N$. Substituting into Equation~\ref{eq:hdi_width_approx}
and setting $w = \omega_{\rm goal}$ gives the exact expected stop iteration:
\begin{equation}\label{eq:expected_stop_iteration_binomial}
N_{\rm goal} = \frac{4 z_{*}^{2}}{\omega_{\rm goal}^{2}}\,\theta(1-\theta) - 1.
\end{equation}
The $-1$ arises from the $N+1$ denominator in the exact Beta posterior variance,
whereas the CLT approximation (Equation~\ref{eq:expected_stop_iteration_clt})
treats the denominator as $N$ and omits it.
As shown in Figure~\ref{fig:min_sample_by_goal} in Section~\ref{sec:expected_stop_iteration},
$N_{\rm goal}$ is maximised at $\theta=0.5$, confirming that the fair coin case requires the
largest sample to achieve a given precision goal, and decreases monotonically as
$\theta$ moves towards $0$ or $1$.
Equation~\ref{eq:pitg_stop_iteration} in the main text uses the CLT approximation
(Equation~\ref{eq:expected_stop_iteration_clt}), which replaces $N+1$ with $N$ and
drops the $-1$ term, negligible for large $N$.
%
% -----------------------------------------------------------------------
% AUTHOR NOTE: Why two prior conventions coexist (Option A decision)
%
% Appendix~\ref{app:hdi} uses the flat prior $\mathrm{Beta}(1,1)$, giving
% posterior $\mathrm{Beta}(s+1, f+1)$, as the principled default for
% posterior computation.  This appendix deliberately retains the Haldane
% prior $\mathrm{Beta}(0,0)$ --- an improper prior (one that does not
% integrate to 1) that assigns infinite weight to the boundary values
% $\theta=0$ and $\theta=1$ --- because it is the only convention that
% yields a clean $-1$ correction in the planning formula.  Here is why
% the two priors lead to different results.
%
% \textbf{Haldane prior} $\mathrm{Beta}(0,0) \;\Rightarrow\;$
% posterior $\mathrm{Beta}(s,\,f)$:
% \begin{equation}
%   \mathrm{Var}(\theta \mid s, f)
%     = \frac{\hat\theta\,(1-\hat\theta)}{N+1}, \qquad \hat\theta = s/N.
% \end{equation}
% Setting $2z_{*}\sqrt{\mathrm{Var}} = \omega_{\rm goal}$ and solving for $N$:
% \begin{equation}
%   N_{\rm goal}
%     = \frac{4z_{*}^{2}}{\omega_{\rm goal}^{2}}\,\theta(1-\theta) - 1.
% \end{equation}
% The $-1$ comes directly from the $N+1$ denominator.
% This is the formula coded in
% \texttt{binomial\_rate\_ci\_width\_to\_sample\_size} (\texttt{utils\_stats.py}):
% \texttt{n\_ = p * (1 - p) / variance\_ - 1}.
%
% \textbf{Flat prior} $\mathrm{Beta}(1,1) \;\Rightarrow\;$
% posterior $\mathrm{Beta}(s+1,\,f+1)$:
% \begin{equation}
%   \mathrm{Var}(\theta \mid s+1, f+1)
%     = \frac{(s+1)(f+1)}{(N+2)^{2}(N+3)}.
% \end{equation}
% This does not simplify to a clean $-1$ correction.
% The leading-order CLT approximation (without any correction term) is the
% practical planning result:
% \begin{equation}
%   N_{\rm goal}
%     \approx \frac{4z_{*}^{2}}{\omega_{\rm goal}^{2}}\,\theta(1-\theta).
% \end{equation}
% Since the CLT approximation already discards $O(1/N)$ terms, and the
% difference between the two priors is also $O(1/N)$, using either
% convention in the planning formula changes the answer by at most 1
% sample for the regime of interest ($N \gg 1$).  The $-1$ correction is
% therefore retained here for consistency with the code, not because it
% is numerically significant.
% -----------------------------------------------------------------------


\subsection*{Continuous Case}

For continuous outcomes when the estimand is a mean $\mu$ with variance $\sigma^{2}$,
two versions of the expected stop iteration are available.

\textbf{Version A (CLT approximation).}
Substituting $V(\mu)=\sigma^{2}$ (estimated by sample variance $s^2$) into
Equation~\ref{eq:expected_stop_iteration_clt} gives
\begin{equation}\label{eq:expected_stop_iteration_mean}
N_{\rm goal} \approx \frac{4 z_{*}^{2}\sigma^{2}}{\omega_{\rm goal}^{2}}.
\end{equation}
This version is accurate for large $N$ where the Student-$t$ critical value
approaches $z_{*}$.

\textbf{Version B (Student-$t$ posterior).}
Under a non-informative prior the posterior for $\mu$ is a scaled Student-$t$
with $\nu = N-1$ degrees of freedom, giving HDI width
$w = 2\,t_{\alpha/2,\,N-1}\,s/\!\sqrt{N}$.
Setting $w = \omega_{\rm goal}$ yields the implicit equation
\begin{equation}\label{eq:expected_stop_iteration_mean_t}
N_{\rm goal} = \frac{4\,t_{\alpha/2,\,N-1}^{2}\,\sigma^{2}}{\omega_{\rm goal}^{2}},
\end{equation}
where $t_{\alpha/2,\,N-1}$ depends on $N$, so the equation must be solved
numerically. For $N \ge 30$ the two versions are nearly identical; Version~A
is preferred for sample-size planning.

